\chapter{DISCUSSION AND EVALUATION}
\label{ch:discussion}

\section{Evaluation against the Requirements}

Table~\ref{tab:reqeval} maps the results of Chapter~\ref{ch:results} to the requirements of Table~\ref{tab:requirements}. A requirement is marked as met only if it was demonstrated on hardware or by an automated test.

\begin{table}[H]
\centering
\small
\caption{Requirement coverage}
\label{tab:reqeval}
\begin{tabularx}{\textwidth}{@{}lXl@{}}
\toprule
\textbf{ID} & \textbf{Evidence} & \textbf{Status} \\
\midrule
R1 & Different kernel clocks and time-quanta counts on the two nodes; error-free exchange at 500\,kbit/s on the physical bus (M4) & met \\
R2 & M1: 100 + 10 frames per second & met \\
R3 & E2E unit tests; M2: 30\,900 frames checked & met \\
R4 & M4: timeout detected within about 100\,ms and out-of-range immediately, output forced to 0\,\%; E2E detection shown with stale frames & met \\
R5 & Sensor ECU button B1 over the bus (M4); Control ECU local fault through the M3 self-test & met \\
R6 & UDS over ISO-TP on 0x7E0/0x7E8 with NRCs: 22 unit tests and 14/14 on-target self-test (loopback); hard reset not yet exercised & met (loopback) \\
R7 & Only MCAL files include the HAL; 34 PC unit tests on the shared modules & met \\
R8 & Hundreds of bus-off events during bus bring-up, each recovered automatically & met \\
\bottomrule
\end{tabularx}
\end{table}

\section{What Loopback Tests Do and Do Not Prove}

Internal loopback validates a large part of the design: the controller configuration, the acceptance filters, the interrupt path, the frame packing, the E2E protection and the application logic. It does not validate the physical layer. In internal loopback the controller acknowledges its own frames, so an error in the transceiver wiring, the termination, or a mismatch in bit timing between the two nodes would stay invisible. The two ECUs deliberately use different kernel clocks and different numbers of time quanta; whether they interoperate is what milestone M4 then showed on the physical bus; a logic-analyser measurement of the bit time and sample point would complete the picture.

\section{Lessons Learned}

\paragraph{Generated code is part of the system under test.} The SysTick defect came from STM32CubeMX output, not from application code. It was only found because the investigation did not stop at ``the tasks do not run'' but checked the symbol table: a weak symbol resolved to \texttt{Default\_Handler} is a concrete, checkable fact. Putting the fix in a user section of the generated file keeps it across regeneration.

\paragraph{On a heterogeneous SoC the M4 inherits the clock tree of Linux.} Code that works on a standalone MCU can fail on the STM32MP1 because peripherals assigned to the M4 still depend on clock settings chosen by Linux. The FDCAN defect was a consequence of a general hardware rule (glitch-free multiplexers need both inputs running) meeting a specific software state (an unused PLL output switched off). The clean long-term solution is to let Linux configure the FDCAN kernel clock through the device tree; the firmware-side switch with read-back is a pragmatic and verified workaround.

\paragraph{Testing the integration, not only the units.} The unit tests showed that the UDS server and the DTC manager behave as specified in isolation. The on-target self-test added what they cannot show: that ISO-TP runs correctly from an RTOS task fed by an interrupt queue, that two tasks can share the CAN transmitter, and that a fault detected by the monitor is visible over UDS with the right status bits. Running the tester on the ECU itself made this possible before a second node exists.

\paragraph{Error counters as a measuring instrument.} Without an oscilloscope, the CAN controller itself was the most informative instrument during the bus bring-up. An acknowledge error proves that a node reads back its own bits, a bit-0 error that the bus does not become dominant, and listen-only mode separates receiving from transmitting. Swapping modules between known-good positions then isolated each segment of the chain in a few minutes.

\paragraph{Debugging without a debugger.} Reading M4 memory and peripheral registers from Linux through \texttt{/dev/mem}, guided by symbol addresses from the ELF file, proved to be an effective technique on this platform, and it also allowed hypotheses to be tested by writing registers (for example enabling PLL4\_R to confirm the multiplexer theory) before changing any firmware.

\paragraph{Evidence before conclusions.} Several intermediate hypotheses during bring-up were wrong: USB cables and adapters were suspected while simply re-seating the connectors resolved the connection problems, the clock calibration unit was suspected but found in bypass, and the role of the HSE kernel-clock gate could not be isolated because its disable register is secure-only. Recording which observations support each conclusion (Section~\ref{sec:bringup}) keeps the report honest about what was actually proven.

\section{Limitations}

\begin{itemize}
  \item The bit timing on the bus has not been measured with a logic analyser; correctness is inferred from error-free operation.
  \item Diagnostics are verified only against a tester on the same core; an external tester on a physical bus is still missing, and the hard reset has not been exercised.
  \item DTCs are kept in RAM only (no non-volatile storage, no aging).
  \item The HAL tick of the M4 runs at about twice the intended rate in production mode.
  \item The FDCAN clock is reconfigured by the M4 rather than by the Linux device tree.
  \item MISRA compliance is checked with the rule subset that cppcheck can decide; there is no certified checker and no manual review of the undecidable rules. The firmware has not been measured for worst-case execution time or interrupt latency.
\end{itemize}
